
\chapter{Discussion}
%------------------------------------------------------------------------------
%%%%%%%%%%%%%%%%%%%%%%%%%%%%%%%%%%%%%%%%%%%%%%%%%%%%%%%%%%%%%%%
%------------------------------------------------------------------------------
\section{Evaluation of the expression of eGFP-tagged uTPs in \textit{Angomonas deanei}}
%------------------------------------------------------------------------------
Understanding the mechanism behind protein targeting of different endosymbiotic systems is crucial to understanding organellogenesis. Thus, in this thesis, the impact of the heterologous expression of the \gls{utp}s inside \glsxtrshort{Adea} was evaluated. The expressed constructs were also attempted to be localized. While the heterologous expression of two \gls{utp}s, although weak, was verified, the localization of \gls{egfp}-tagged \gls{utp}s was not possible. 

In conclusion, it was not possible to demonstrate that the nitroplast transit peptide affected or localized within \glsxtrshort{Adea}; therefore, the targeting mechanisms for both endosymbiotic systems remained elusive. 
%------------------------------------------------------------------------------
\subsection{Expression levels}
The integration of the transfection cassettes containing the \gls{egfp}-tagged \gls{utp}s was achieved for only two out of the three constructs generated, with \gls{utp_C} missing from the analyses. The Western blot analysis and the in-gel fluorescence assay verified that \glsxtrshort{Adea} was able to express \gls{egfp}-\gls{utp_B} and \gls{egfp}-\gls{utp_E} at similar levels under normal growth conditions. Moreover, the strains expressing the \gls{egfp}-\gls{utp} constructs showed no phenotypic deviation from the \gls{WT} when examined under the microscope. Hence, it is likely that the nitroplast transit peptides do not interact negatively with the cellular processes or structures in \glsxtrshort{Adea}. With no toxic effects present, it should be possible to integrate and express the transfection cassette containing \gls{utp_C}, giving similar results as the other two \gls{utp}s.   

\par Though the \gls{utp}s may not be toxic to \glsxtrshort{Adea}, there was a clear difference in the expression levels of \gls{egfp} by itself and fused to the nitroplast targeting signals. The expression of \gls{egfp} by itself from the \gls{g-ama} locus displayed higher expression levels than compared to the \gls{egfp}-\gls{utp} constructs expressed from the \gls{d-ama} locus. It seems unlikely that the observed difference in expression could be attributed to the effect of the different loci, as the transcriptome data showed that the \gls{d-ama} locus ranked among the top 40 and the \gls{g-ama} locus close to the top 100 most transcribed genes in \glsxtrshort{Adea} \parencite{morales_development_2016}. Moreover, the possibility of degradation of the fusion proteins due to being misfolded would be unlikely since gls{egfp}-\gls{utp} constructs showed fluorescence, indicating correct folding for at least \gls{egfp} and potentially also for the \gls{utp}. Additionally, the export of the recombinant transit peptides 
to the outside of the cells did not occur, as it was verified by the Western blot analysis of the supernatant samples from the expression cultures, underlining that it was indeed a difference in expression levels. However, a possible reason for the different expression levels might be that the codon usage bias of \glsxtrshort{Adea} differed from the codon usage of the nucleotide sequence used for the \gls{utp}s, therefore decreasing expression, which had been demonstrated for other trypanosomatids \parencite{de_freitas_nascimento_codon_2018, horn_codon_2008}. Though speculative, it might be that more secondary structures in the mRNA formed because of non-favorable codon usage combined with the repetitiveness of the nucleotide sequence due to using one codon per amino acid \parencite{subramanian_comparison_2015}.

\subsection{Localization}
The localization of the \gls{egfp}-tagged \gls{utp}s was not successful despite the verified heterologous expression. From the Western blot analysis of the supernatant of the harvested culture, which lacked signals for the strains expressing the \gls{egfp}-\gls{utp} constructs, it could be concluded that the \gls{egfp}-tagged nitroplast transit peptides must localize inside \glsxtrshort{Adea}. The strains expressing these fusion proteins were examined using two different microscopy methods, epifluorescence and confocal. Neither method detected \gls{egfp} signals for the \gls{utp} constructs. In contrast, the detection via the in-gel fluorescence assay, using samples from harvested cultures, was possible. The issue most likely was the levels at which the \gls{egfp}-tagged \gls{utp}s were present inside the cells, since \gls{egfp} folded correctly even attached to the nitroplast transit peptide. One possible explanation for the lack of \gls{egfp} signals, when taking the weak heterologous expression of these constructs into consideration, could be that the \gls{egfp}-tagged nitroplast targeting signals diffused throughout the cytosole, and being present at such low concentrations would require more sensitive microscopy, like single-molecule localization microscopy \parencite{lelek_single-molecule_2021}. An alternative explanation could be that the \gls{egfp} signals from the \gls{utp} constructs were outside the focus plane, since only the \gls{ES} was in focus, henceforth not being detected. This would also assume that the \gls{utp}s did not diffuse in the cytosole but instead localized to specific regions within the cells, for which we have no evidence so far
%------------------------------------------------------------------------------
%%%%%%%%%%%%%%%%%%%%%%%%%%%%%%%%%%%%%%%%%%%%%%%%%%%%%%%%%%%%%%%
%------------------------------------------------------------------------------
\section{Possible purification of 6xHis-tagged uTPs}
%------------------------------------------------------------------------------
To gain further insights into the mechanisms behind protein targeting and import into the nitroplast of \glsxtrshort{Bbig}, it is necessary to study the structures of the corresponding transit peptides. Thus, in the scope of this thesis, the expression conditions and test purification of two \gls{6xhis}-tagged \gls{utp} constructs were analyzed, providing groundwork for larger-scale expression and purification runs.  
%------------------------------------------------------------------------------
\subsection{Expression and solubility}
Expression of both \gls{6xhis}-tagged \gls{utp}s was possible across different expression conditions. Almost all the cultures grown at 16~°C, except EC174 containing the \gls{utp_E} grown in \gls{2yt}, did not express the recombinant transit peptides. Additionally, the overnight samples compared to the 4~hours after induction samples yielded less \gls{utp}s. These two facts suggested that either degradation occurred or that the \gls{utp}s formed inclusion bodies.  

The solubility test verified expression of soluble \gls{6xhis}-\gls{utp}s within some of the selected expression conditions. The most soluble proteins were obtained when cultures were grown at 28~°C. Expression cultures grown at temperatures above or below 28~°C had decreased amounts of soluble fusion proteins, likely due to the formation of inclusion bodies. However, based on the finding that the \gls{utp}s likely contain transmembrane helices in the C-terminal region, it could be postulated that the recombinant transit peptides inserted themselves into the membrane of \glsxtrshort{Eco}, therefore not being present in the cytosolic fraction of the cells. That would require the \gls{utp}s to interact with bacterial cell membranes, which had not been demonstrated yet.   
%Further, it could be stated that the \gls{utp}s exhibited similar behavior when it came to expression, indicating consistent characteristics among them. Thus, it would be likely that utp_C has similar expression conditions.
%------------------------------------------------------------------------------
%------------------------------------------------------------------------------
\subsection{Outcomes of the test purification}
Despite the low solubility, the \gls{6xhis}-tagged \gls{utp}s were able to be purified at typical concentrations of around 250~--500~mM imidazole. Though most of the samples appeared to remain in the flow-through, suggesting that they did not bind to the resin. This likely arose from the methodology used for purification, such that too much protein was loaded onto the resin. For this reason, it could be assumed that this issue will be resolved when \gls{fplc} is performed using the Äkta system. 

Yet, what remained unclear was the identity of the smaller protein bands, eluting earlier than the \gls{utp}s (50~--~100~mM imidazole). Endogenous contamination could be excluded since the LOBSTR strain was specifically engineered to avoid such occurrences \parencite{andersen_optimized_2013}. It could be speculated that these bands contained truncated \gls{6xhis}-\gls{utp}s, though there was no evidence of cleavage by proteases for the \gls{utp}s. 
    
%------------------------------------------------------------------------------
%%%%%%%%%%%%%%%%%%%%%%%%%%%%%%%%%%%%%%%%%%%%%%%%%%%%%%%%%%%%%%%
%------------------------------------------------------------------------------
\section{Similarities and differences to other transit peptides}
To gain further knowledge about the structure of the \gls{utp}s and a possible mechanism of protein transport, bioinformatic analyses were conducted. These analyses provided evidence for the potential structure as well as revealed similarities between the \gls{utp}s and other transit peptides. Interestingly, nitroplast-targeted proteins might be transported to the organelle in a similar way that chromatophore-targeted proteins are.   
%------------------------------------------------------------------------------
\subsection{Structure of UCYN-A transit peptides}
The \gls{utp}s are comparable in size to \gls{crtps}, which are about 200 amino acids long \parencite{singer_massive_2017}. Whereas \gls{mtps} and \gls{ctps} are much smaller (15~--~80 amino acids) \parencite{garg_role_2016, von_heinje_chloroplast_1991}. An increase in the length of the transit would essentially lead to an increase in its structural complexity. While \gls{ctps} tend to form random coils in aqueous and α-helices in hydrophobic environments, \gls{crtps} display two structured domains \parencite{oberleitner_bipartite_2022, klimenko_paulinella_2025, bruce_chloroplast_2000}. Similarly, the \gls{utp}s were predicted to form two rather distinct domains, the N- and C-terminal regions. The N-terminal region contained α-helical structures that were arranged in a U-shape, with potentially two helices running antiparallel. This seemed to be consistent across the different \gls{utp}s. This structure was also predicted by a team of student researchers who constructed two concatenated consensus sequences that were predicted to fold the same \parencite{purrinos_nitroblast_2024}. Though hypothetical, this structure would be similar to intra-chain coiled coils, which have been described as molecular recognition modules and intramolecular spacers \parencite{hartmann_functional_2017}. Any of these two functions could be possible for the \gls{utp}s because the \gls{utp}s would either (i) separate the cargo protein from the C-terminal region for proper recognition of a putative receptor or (ii) act as the recognition site for a putative receptor. However, PHYRE did not identify direct structural homologs for the N-terminal regions of the \gls{utp}s and instead gave multiple low-confidence hits to α-helical transmembrane domains. Moreover, the C-terminal region was also predicted to contain α-helices. The confidence level for this prediction was rather low compared to the N-terminal regions. This might be because Alphafold handles very flexible amino acid stretches or intrinsically disordered domains less effectively \parencite{ruff_alphafold_2021}. This in turn could suggest that the C-terminal region of the \gls{utp}s might be disordered or at least flexible. Additionally, TargetP and PHYRE also assigned transmembrane domains to the C-terminal region. Likewise, a transmembrane α-helix had been described for part 1 of the \gls{crtps} that is suggested to be inserted into the membrane of the endoplasmic reticulum \parencite{singer_massive_2017}. Furthermore, based on the observation that these \gls{utp}s occur at the C-terminus of the cargo proteins, it could be concluded that the targeted proteins would be transported in a folded state. That would be in contrast to how proteins targeted to the mitochondria/chloroplasts are transported to their respective organelles \parencite{schleiff_common_2011}. Though this might be more in line with how proteins are transported to the chromatophore, possibly sharing similar mechanisms \parencite{oberleitner_bipartite_2022, singer_massive_2017}. Alternatively, it was suggested that the bipartite structure of the \gls{utp}s was rather for sub-organellar localization  \parencite{purrinos_nitroblast_2024}.

%------------------------------------------------------------------------------
%------------------------------------------------------------------------------
\subsection{Physicochemical properties}
The alignment of the \gls{utp}s revealed that they displayed a similar amount of diversity as it was described for other transit peptides \parencite{lee_understanding_2021, huang_refining_2009}. Despite such diversity, some sites were conserved, which was in accordance with the concatenated motif sequence created by that student research team \parencite{purrinos_nitroblast_2024}. The motif sequence that they generated seemed to exclude the C-terminal region, even though the \gls{utp}s studied here did display potential conserved sites. Moreover, the predicted α-helices in the N-terminal region showed amphipathicity as analyzed by Heliquest. These amphipathic helices seemed most similar to \gls{ctps} because the polar face lacked any specific charge \parencite{bruce_chloroplast_2000}. In contrast, the C-terminal region was more hydrophobic, having minimal hydrophobic moments, which would fit its prediction as a transmembrane domain by TargetP and PHYRE. 

Overall, it could be postulated that the physicochemical properties of the \gls{utp}s are more important than their primary amino acid sequence, similar to what is observed in \gls{mtps} \gls{ctps}. Further, the transport of these proteins containing the \gls{utp}s might occur post-translationally in a folded state via the secretory pathway, like the proteins being targeted to the chromatophore. 
%------------------------------------------------------------------------------
%%%%%%%%%%%%%%%%%%%%%%%%%%%%%%%%%%%%%%%%%%%%%%%%%%%%%%%%%%%%%%%
%------------------------------------------------------------------------------
\section{Outlook}
%------------------------------------------------------------------------------
%------------------------------------------------------------------------------
\subsection{Heterologous expression and localization within endosymbiotic systems}
The heterologous expression of \gls{egfp}-tagged \gls{utp}s was possible but could be optimized for higher expression levels. For that purpose, integrating and simultaneously expressing the recombinant \gls{utp}s into different loci might resolve this issue. Furthermore, an immunofluorescence assay with an \gls{egfp} booster could be performed to enhance the signal strength of the \gls{egfp}-\gls{utp}s. Additionally, confocal fluorescence microscopy with focus stacking should be attempted next for localization studies to not miss out on any potential signals of the \gls{egfp}+\gls{utp} constructs. With these adjustments, it might then be possible to localize \gls{utp}s and answer the question if the \gls{utp}s interact with different endosymbiotic systems.
%------------------------------------------------------------------------------
%------------------------------------------------------------------------------
\subsection{Next steps in the purification process of UCYN-A transit peptides}
Following the successful test runs, larger cultures could be set up for bulk expression, growing the cultures at 28~°C 4~hours after induction in \gls{lb} for \gls{utp_C} and \gls{2yt} for \gls{utp_E}. A test run using the \gls{fplc} could then be conducted. Based on the results of the conducted tests in this thesis, it seemed unlikely that more optimization would be required for purification. Though salt concentrations of the buffers, the metal ion used for binding \gls{6xhis}, and the presence of imidazole in the washing buffer could all be adjusted to maximise purified \gls{utp}s. Following \gls{fplc}, it would be advised to perform size exclusion chromatography to avoid co-purification of the smaller protein band. With all that, the groundwork for obtaining the structure for the nitroplast targeting signals was laid out.
%------------------------------------------------------------------------------
%------------------------------------------------------------------------------
\subsection{Bioinformatic analyses}
Through the bioinformatic analyses, it was possible to gain a deeper understanding of the \gls{utp}s. The verification of the predicted transmembrane domains could be done by running an \gls{page} with samples from differential centrifugation fractions of the \glsxtrshort{Eco} solubility samples treated with different detergents. Moreover, Foldseek could be utilized to obtain models with domains similar to those of the \gls{utp}s that would also have a higher confidence level than models PHYRE provided \parencite{van_kempen_fast_2024}.

\section{Final remarks}
Studying the nitroplast targeting signals to discover their underlying principles is not only interesting for understanding the establishment of organelles, but also might lead to the achievement of synthetic nitroplast in plants in the future, which is of high economic interest as it would eliminate the need for fertilizer with added N\textsubscript{2}. Though the research on these new transit peptides has just begun, promising results are to be expected